\documentclass[11pt]{article}
\usepackage[margin=1in]{geometry}
\usepackage{enumitem}
% =============================================================================
% Preambles/header.tex — shared LaTeX/Beamer preamble for this project
%
% Use from a Beamer lecture by prepending:
%     \input{header}
% and compiling with TEXINPUTS=../Preambles:$TEXINPUTS (the /compile-latex
% skill does this automatically).
%
% PALETTE CONTRACT. Color names defined here MUST match the names in
% Quarto/theme-template.scss so Beamer and Quarto renderings use the same
% palette. The scripts/check-palette-sync.sh script enforces this.
%
% To customize: edit the HEX values below AND the matching $variable in
% Quarto/theme-template.scss, then run scripts/check-palette-sync.sh.
% =============================================================================

% -----------------------------------------------------------------------------
% Engine detection + encoding
%
% Our /compile-latex skill uses XeLaTeX, where inputenc is unnecessary and
% can produce encoding warnings. Only load inputenc under pdfTeX.
% -----------------------------------------------------------------------------
\usepackage{iftex}
\ifPDFTeX
  \usepackage[utf8]{inputenc}
\fi

\usepackage{amsmath, amssymb, amsthm}
\usepackage{booktabs}
\usepackage{graphicx}

% -----------------------------------------------------------------------------
% PALETTE — mirrors Quarto/theme-template.scss variable names.
% Load xcolor and define colors BEFORE anything that references them
% (notably hyperref, hypersetup, and the Beamer color assignments).
% -----------------------------------------------------------------------------
\usepackage{xcolor}

\definecolor{primary-blue}{HTML}{012169}      % headings, accents
\definecolor{primary-gold}{HTML}{B9975B}      % emphasis, borders
\definecolor{highlight-yellow}{HTML}{F2A900}  % markers, alerts
\definecolor{light-bg}{HTML}{E8EDF5}          % title / section backgrounds
\definecolor{jet}{HTML}{1A1A1A}               % body text

% Semantic colors (used in diagrams and annotations)
\definecolor{positive}{HTML}{15803D}          % good / observed / identified
\definecolor{negative}{HTML}{B91C1C}          % bad / problematic / bias
\definecolor{neutral}{HTML}{525252}           % reference / context

% Highlight accents (match .hi-* classes in the SCSS)
\definecolor{hi-slate}{HTML}{314F4F}
\definecolor{hi-green}{HTML}{2E8B57}
\definecolor{hi-red}{HTML}{C41E3A}

% Semantic aliases so skill templates and snippets can write
% \textcolor{accent}{...} without hard-coding "primary-blue".
\colorlet{accent}{primary-blue}
\colorlet{accent2}{primary-gold}

% -----------------------------------------------------------------------------
% Hyperref — loaded AFTER xcolor + palette so link colors resolve.
% -----------------------------------------------------------------------------
\usepackage{hyperref}
\hypersetup{colorlinks=true, linkcolor=primary-blue, urlcolor=primary-blue,
            citecolor=primary-blue, pdfborder={0 0 0}}

% -----------------------------------------------------------------------------
% Course code — set once per deck (after \input{header}, before \begin{document})
% via \coursecode{CS401}. Shown in the footer on every slide so a deck's course
% is identifiable without opening the title page. Empty by default.
% -----------------------------------------------------------------------------
\makeatletter
\newcommand{\coursecode}[1]{\gdef\@coursecodetext{#1}}
\gdef\@coursecodetext{}
\makeatother

% -----------------------------------------------------------------------------
% Beamer theme assignments (only apply under Beamer).
%
% We detect Beamer via \@ifclassloaded{beamer}, which is the documented,
% non-internal way. This must be wrapped in \makeatletter/\makeatother
% because \@ifclassloaded contains an @-catcode character.
% -----------------------------------------------------------------------------
\makeatletter
\@ifclassloaded{beamer}{%
  \usefonttheme{professionalfonts}
  \setbeamercolor{structure}{fg=primary-blue}
  \setbeamercolor{frametitle}{fg=primary-blue}
  \setbeamercolor{title}{fg=primary-blue}
  \setbeamercolor{subtitle}{fg=primary-gold}
  \setbeamercolor{author}{fg=primary-blue}
  \setbeamercolor{institute}{fg=neutral}
  \setbeamercolor{date}{fg=primary-gold}
  \setbeamercolor{itemize item}{fg=primary-blue}
  \setbeamercolor{itemize subitem}{fg=primary-gold}
  \setbeamercolor{alerted text}{fg=primary-gold}
  \setbeamercolor{block title}{fg=white, bg=primary-blue}
  \setbeamercolor{block body}{bg=light-bg}
  % Semantic block variants: block=definition/concept (blue), exampleblock=
  % worked example (green/positive), alertblock=key takeaway or caution (gold).
  % Keep to <=2 colored boxes per slide across ALL three variants (INV-7).
  \setbeamercolor{block title example}{fg=white, bg=positive}
  \setbeamercolor{block body example}{bg=light-bg}
  \setbeamercolor{block title alerted}{fg=white, bg=primary-gold}
  \setbeamercolor{block body alerted}{bg=light-bg}

  % Minimal footer: course code (left) + page X / N (right), no navigation clutter
  \setbeamertemplate{navigation symbols}{}
  \setbeamertemplate{footline}{%
    \color{neutral}\footnotesize\hspace{0.7em}\@coursecodetext\hfill
    \insertframenumber/\inserttotalframenumber\hspace{0.7em}\vspace{0.3em}}
}{}
\makeatother

% -----------------------------------------------------------------------------
% TikZ setup — libraries the snippet gallery uses
% -----------------------------------------------------------------------------
\usepackage{tikz}
\usetikzlibrary{arrows.meta, positioning, calc, decorations.pathreplacing,
                fit, shapes.geometric, backgrounds}

% Shared TikZ styles — snippets and hand-written diagrams can pick these up
% via `\begin{tikzpicture}[every node/.style={...}]` or by naming specific
% styles (e.g., `\node[dag-node] (X) at (0,0) {X};`).
\tikzset{
  % Node styles for causal DAGs and similar
  dag-node/.style = {
    draw=primary-blue, fill=light-bg, rounded corners=2pt,
    minimum width=1.2cm, minimum height=0.9cm, align=center, font=\small
  },
  decision-node/.style = {
    draw=primary-gold, fill=white, diamond, aspect=1.8,
    minimum width=1.8cm, minimum height=1.2cm, align=center, font=\small,
    inner sep=1pt
  },
  % Edge styles — observed vs counterfactual
  observed-edge/.style = {-{Latex[length=2.5mm]}, thick, draw=jet},
  counterfactual-edge/.style = {-{Latex[length=2.5mm]}, thick, draw=neutral, dashed},
  confound-edge/.style = {-{Latex[length=2.5mm]}, thick, draw=negative, dashed},
  % Data-point styles for plots
  observed-dot/.style = {fill=primary-blue, draw=primary-blue, circle, inner sep=1.2pt},
  counterfactual-dot/.style = {fill=white, draw=neutral, circle, inner sep=1.2pt}
}

% -----------------------------------------------------------------------------
% Convenience macros
% -----------------------------------------------------------------------------
\newcommand{\muted}[1]{\textcolor{neutral}{#1}}
\newcommand{\key}[1]{\textcolor{primary-gold}{\textbf{#1}}}
\newcommand{\good}[1]{\textcolor{positive}{#1}}
\newcommand{\bad}[1]{\textcolor{negative}{#1}}

% Standout frame macro: full-bleed dark background for section transitions.
% Beamer-only (uses \begin{frame}/\setbeamercolor/\usebeamerfont) -- do not
% call from an article-class file (e.g. Notes/<CODE>/*-notes.tex). \muted,
% \key, \good, \bad, and \coursecode above are plain xcolor/gdef and are
% safe to use from either class; verified when Notes/article-class support
% was added.
% Expand: \transitionslide{Your transition title}
\newcommand{\transitionslide}[1]{%
  \begingroup
  \setbeamercolor{background canvas}{bg=primary-blue}
  \begin{frame}[plain]
    \centering
    \vfill
    {\usebeamerfont{title}\color{white}\Large #1\par}
    \vfill
  \end{frame}
  \endgroup
}

% Section divider frame (Beamer-only), an alternative to \transitionslide with a
% darker two-line style: near-black (jet) background, a small white label line
% above a larger gold title, and a thin gold rule along the bottom edge.
% Expand: \sectiondivider{Section 2}{Pointers}
\newcommand{\sectiondivider}[2]{%
  \begingroup
  \setbeamercolor{background canvas}{bg=jet}
  \begin{frame}[plain]
    \vspace*{3.0cm}
    {\color{white}\ttfamily\large #1\par}
    \vspace{0.5cm}
    {\color{primary-gold}\LARGE #2\par}
    \begin{tikzpicture}[remember picture, overlay]
      \draw[primary-gold, line width=2.5pt]
        ($(current page.south west)+(0,0.1)$) -- ($(current page.south east)+(0,0.1)$);
    \end{tikzpicture}
  \end{frame}
  \endgroup
}

% -----------------------------------------------------------------------------
% Channel watermark — "Concepts That Click" (YouTube).
%
% Article-class files (CompetitiveExam Q/A sets, Notes, Assignments) get a
% light diagonal draft watermark on every page plus a centered footer naming
% the channel. Beamer decks get a subtle TikZ background watermark instead
% (the footline already carries the course code). Centralized here so every
% MCQ PDF is branded without per-file edits.
% -----------------------------------------------------------------------------
\makeatletter
\@ifclassloaded{article}{%
  \usepackage{draftwatermark}
  \SetWatermarkText{Concepts That Click}
  \SetWatermarkScale{1}
  \SetWatermarkFontSize{2cm}
  \SetWatermarkLightness{0.86}
  \SetWatermarkAngle{45}
  \usepackage{fancyhdr}
  \pagestyle{fancy}
  \fancyhf{}
  \fancyfoot[C]{\footnotesize\textcolor{neutral}{Concepts That Click~|~YouTube: \href{https://www.youtube.com/@concepts.thatclick}{@concepts.thatclick}}}
  \renewcommand{\headrulewidth}{0pt}
  \renewcommand{\footrulewidth}{0pt}
  \fancypagestyle{plain}{%
    \fancyhf{}%
    \fancyfoot[C]{\footnotesize\textcolor{neutral}{Concepts That Click~|~YouTube: \href{https://www.youtube.com/@concepts.thatclick}{@concepts.thatclick}}}%
    \renewcommand{\headrulewidth}{0pt}%
    \renewcommand{\footrulewidth}{0pt}%
  }
}{}
\@ifclassloaded{beamer}{%
  \setbeamertemplate{background}{%
    \begin{tikzpicture}[remember picture, overlay]
      \node[rotate=45, opacity=0.055, align=center]
        at (current page.center)
        {\Huge\textcolor{neutral}{Concepts That Click}};
    \end{tikzpicture}%
  }
}{}
\makeatother 
\coursecode{JKSSB}
\title{Lesson 35 Practice: URL, DNS, HTTP and HTTPS}
\author{JKSSB Computer Awareness}
\date{\today}
\begin{document}
\maketitle
\noindent\textbf{Provenance:} All 20 questions are original JKSSB-pattern
questions (\textit{[Original]}); none reproduces an official paper item.

\begin{enumerate}[label=\textbf{Q\arabic*.}, leftmargin=*]
\item \textit{[Original]} Expand URL.
  \begin{enumerate}[label=\Alph*.]
    \item Uniform Resource Locator
    \item Universal Record Link
    \item Unified Reference Locator
    \item Uniform Retrieval Language
  \end{enumerate}
\item \textit{[Original]} Expand DNS.
  \begin{enumerate}[label=\Alph*.]
    \item Domain Name System
    \item Data Network Service
    \item Digital Node Scheme
    \item Distributed Number Sequence
  \end{enumerate}
\item \textit{[Original]} Expand HTTPS.
  \begin{enumerate}[label=\Alph*.]
    \item Hypertext Transfer Protocol Secure
    \item High Transfer Text Protocol System
    \item Hyperlink Text Transmission Protocol Standard
    \item Host Transfer Protocol Service
  \end{enumerate}
\item \textit{[Original]} In the URL
  \texttt{https://www.example.in:443/study/computer?topic=dns}, which part
  is the scheme?
  \begin{enumerate}[label=\Alph*.]
    \item \texttt{https}
    \item \texttt{www.example.in}
    \item \texttt{:443}
    \item \texttt{/study/computer}
  \end{enumerate}
\item \textit{[Original]} In \texttt{mail.example.com}, which label is the
  subdomain?
  \begin{enumerate}[label=\Alph*.]
    \item \texttt{.com}
    \item \texttt{example}
    \item \texttt{mail}
    \item \texttt{mailto}
  \end{enumerate}
\item \textit{[Original]} Which of the following is a Top-Level Domain
  (TLD)?
  \begin{enumerate}[label=\Alph*.]
    \item \texttt{www}
    \item \texttt{.in}
    \item \texttt{/study}
    \item \texttt{?topic=dns}
  \end{enumerate}
\item \textit{[Original]} The default port for HTTP is
  \underline{\hspace{1.2cm}}, and for HTTPS it is
  \underline{\hspace{1.2cm}}.
  \begin{enumerate}[label=\Alph*.]
    \item 443; 80
    \item 80; 443
    \item 25; 110
    \item 80; 80
  \end{enumerate}
\item \textit{[Original]} Which part of a URL carries extra parameters such
  as \texttt{?topic=dns\&lang=en}?
  \begin{enumerate}[label=\Alph*.]
    \item Scheme
    \item Domain name
    \item Path
    \item Query
  \end{enumerate}
\item \textit{[Original]} Which statement correctly distinguishes a URL from
  a domain name?
  \begin{enumerate}[label=\Alph*.]
    \item A URL is the full address of one resource; the domain name is only
      the host part inside it.
    \item A domain name is the full address of one resource; the URL is only
      the host part inside it.
    \item A URL and a domain name are identical in every part.
    \item A URL names the server hardware; a domain name names the page
      content.
  \end{enumerate}
\item \textit{[Original]} IPv4 is a \underline{\hspace{1.2cm}} address,
  whereas IPv6 is a \underline{\hspace{1.2cm}} address.
  \begin{enumerate}[label=\Alph*.]
    \item 128-bit; 32-bit
    \item 32-bit; 128-bit
    \item 64-bit; 256-bit
    \item 32-bit; 256-bit
  \end{enumerate}
\item \textit{[Original]} The main function of DNS is to:
  \begin{enumerate}[label=\Alph*.]
    \item translate domain names into IP addresses.
    \item encrypt web pages during transfer.
    \item assign a port number to every browser.
    \item carry e-mail between mail servers.
  \end{enumerate}
\item \textit{[Original]} Which of the following is \textbf{NOT} true of DNS?
  \begin{enumerate}[label=\Alph*.]
    \item It maps domain names to IP addresses.
    \item It uses a distributed hierarchy of servers.
    \item It caches answers to speed up repeat visits.
    \item It encrypts the contents of web pages.
  \end{enumerate}
\item \textit{[Original]} When a domain name is not cached, the resolver
  obtains its IP address by:
  \begin{enumerate}[label=\Alph*.]
    \item asking the DNS hierarchy (root, then TLD, then authoritative
      server) for the address.
    \item guessing the address from the URL path.
    \item reading the page content over HTTP first.
    \item copying the port number into the address field.
  \end{enumerate}
\item \textit{[Original]} HTTP stands for:
  \begin{enumerate}[label=\Alph*.]
    \item Hypertext Transfer Protocol
    \item Hyperlink Transmission Process
    \item Host Text Transfer Program
    \item High-speed Transport Protocol
  \end{enumerate}
\item \textit{[Original]} TLS stands for:
  \begin{enumerate}[label=\Alph*.]
    \item Transport Layer Security
    \item Transfer Link Service
    \item Terminal Login System
    \item Transmission Line Splitter
  \end{enumerate}
\item \textit{[Original]} Which statement about HTTP and HTTPS is correct?
  \begin{enumerate}[label=\Alph*.]
    \item HTTP encrypts pages by default; HTTPS removes that encryption.
    \item HTTPS is HTTP protected by TLS; plain HTTP provides no encryption.
    \item HTTP uses port 443 and HTTPS uses port 80.
    \item HTTPS is a page-design language unrelated to HTTP.
  \end{enumerate}
\item \textit{[Original]} A browser shows a padlock for
  \texttt{https://www.example.in} because:
  \begin{enumerate}[label=\Alph*.]
    \item the TLS handshake authenticated the server and encrypted the
      channel.
    \item DNS refused to resolve the name.
    \item the page was transferred over port 80.
    \item the query string was removed from the URL.
  \end{enumerate}
\item \textit{[Original]} Which sequence gives the correct flow when a user
  types a secure URL and presses Enter?
  \begin{enumerate}[label=\Alph*.]
    \item URL $\rightarrow$ DNS $\rightarrow$ IP $\rightarrow$ server
      $\rightarrow$ TLS handshake $\rightarrow$ page.
    \item IP $\rightarrow$ URL $\rightarrow$ DNS $\rightarrow$ TLS handshake
      $\rightarrow$ server $\rightarrow$ page.
    \item URL $\rightarrow$ IP $\rightarrow$ page $\rightarrow$ DNS
      $\rightarrow$ server $\rightarrow$ TLS handshake.
    \item DNS $\rightarrow$ TLS handshake $\rightarrow$ URL $\rightarrow$ IP
      $\rightarrow$ page $\rightarrow$ server.
  \end{enumerate}
\item \textit{[Original]} Evaluate the statements:
  \begin{enumerate}[label=\Roman*.]
    \item DNS translates domain names into IP addresses.
    \item IPv6 is a 128-bit address.
    \item Plain HTTP encrypts the data it carries.
  \end{enumerate}
  \begin{enumerate}[label=\Alph*.]
    \item I only
    \item I and II only
    \item II and III only
    \item I, II, and III
  \end{enumerate}
\item \textit{[Original]} Which pairing is correctly matched?
  \begin{enumerate}[label=\Alph*.]
    \item Scheme --- \texttt{www.example.in}
    \item Port 80 --- default for HTTPS
    \item Query --- parameters after \texttt{?} in a URL
    \item IPv6 --- 256-bit address
  \end{enumerate}
\end{enumerate}
\end{document}
